% ============================================================
\chapter*{Preface}
\addcontentsline{toc}{chapter}{Preface}
\markboth{Preface}{Preface}
% ============================================================

Topological Data Analysis (TDA) is a relatively young field, born in the
early 2000s at the confluence of algebraic topology, computational geometry,
and data science. Its fundamental goal is to extract \emph{qualitative} and
\emph{structural} information from complex datasets by exploiting topological
invariants that are robust to noise and continuous deformations.

\medskip

While classical statistical methods rely on parametric models or distributional
assumptions, TDA adopts a radically different perspective: it seeks to
characterize the \emph{shape} of data. This approach has proven particularly
fruitful in domains as diverse as molecular biology, neuroscience, materials
science, quantitative finance, and deep learning.

%% ----------------------------------------------------------------
\section*{Course Objectives}

This course is designed for Master's and PhD students with a background in:
\begin{itemize}
  \item \textbf{Algebra and topology}: groups, rings, modules; topological
        spaces, compactness, connectedness.
  \item \textbf{Analysis and probability}: metric spaces, measure theory,
        convergence.
  \item \textbf{Algorithms and programming}: complexity, data structures,
        Python.
\end{itemize}

Upon completion, the student will be able to:
\begin{enumerate}
  \item Construct and manipulate the main simplicial complexes
        (Vietoris--Rips, \v{C}ech, Alpha) associated with point clouds.
  \item Compute simplicial homology and persistent homology, both
        theoretically and algorithmically.
  \item Interpret persistence diagrams and barcodes, and assess their
        stability.
  \item Apply the Mapper algorithm for exploratory visualization.
  \item Integrate topological descriptors into machine learning and deep
        learning pipelines.
  \item Conduct a research project applying TDA to real-world data.
\end{enumerate}

%% ----------------------------------------------------------------
\section*{Course Organization}

The course is structured into eleven chapters, grouped into four parts:

\begin{description}
  \item[Part~I --- Foundations (Chapters~1--3)]
    Motivations, simplicial complexes, homology, and introduction to
    persistent homology.
  \item[Part~II --- Persistence Theory (Chapters~4--6)]
    Persistence diagrams, stability theorems, and construction of
    filtered complexes.
  \item[Part~III --- Methods and Algorithms (Chapters~7--8)]
    Mapper algorithm, distances between persistence diagrams
    (Wasserstein, bottleneck).
  \item[Part~IV --- TDA and Learning (Chapters~9--11)]
    Integration with machine learning, concrete applications, and
    connections to deep learning.
\end{description}

Each chapter contains rigorous definitions, theorems with proofs (or proof
sketches), detailed examples, Python implementations using \texttt{gudhi},
\texttt{ripser}, and \texttt{giotto-tda}, as well as exercises of
increasing difficulty.

%% ----------------------------------------------------------------
\section*{Notation and Conventions}

\begin{center}
\renewcommand{\arraystretch}{1.3}
\begin{tabular}{cl}
\toprule
\textbf{Symbol} & \textbf{Meaning} \\
\midrule
$\R^d$ & Euclidean space of dimension $d$ \\
$\N, \Z, \Q$ & Natural numbers, integers, rationals \\
$\mathbb{F}$ & Finite field (often $\mathbb{F}_2 = \Z/2\Z$) \\
$K$ & Abstract simplicial complex \\
$\sigma = [v_0, \ldots, v_p]$ & $p$-simplex \\
$C_p(K)$ & Group of $p$-chains of $K$ \\
$\partial_p$ & Boundary operator in dimension $p$ \\
$H_p(K)$ & $p$-th homology group of $K$ \\
$\beta_p$ & $p$-th Betti number ($= \dim H_p$) \\
$\mathrm{VR}(X, r)$ & Vietoris--Rips complex at parameter $r$ \\
$\check{C}(X, r)$ & \v{C}ech complex at parameter $r$ \\
$\mathrm{Dgm}_p$ & Persistence diagram in dimension $p$ \\
$d_B, W_p$ & Bottleneck distance, $p$-Wasserstein distance \\
$\norm{\cdot}$ & Norm (Euclidean by default) \\
$\abs{S}$ & Cardinality of set $S$ \\
$\inner{\cdot}{\cdot}$ & Inner product \\
\bottomrule
\end{tabular}
\end{center}

%% ----------------------------------------------------------------
\section*{Software Environment}

The implementations in this course use \textbf{Python~3.10+} with the
following libraries:

\begin{itemize}
  \item \textbf{GUDHI} (\url{https://gudhi.inria.fr}): C++/Python library
        developed at INRIA, offering a wide range of simplicial structures,
        persistence algorithms, and statistical tools.
  \item \textbf{Ripser} (\url{https://ripser.scikit-tda.org}): ultra-fast
        implementation of persistent homology for Vietoris--Rips complexes.
  \item \textbf{giotto-tda} (\url{https://giotto-ai.github.io/gtda-docs}):
        machine-learning oriented library, scikit-learn compatible.
  \item \textbf{scikit-learn}, \textbf{NumPy}, \textbf{SciPy},
        \textbf{Matplotlib}: preprocessing, numerical computation, and
        visualization.
  \item \textbf{PyTorch}: for deep learning models integrating topological
        layers.
\end{itemize}

\begin{warning}[title=\textbf{Recommended Installation}]
We recommend using a virtual environment:
\begin{verbatim}
python -m venv tda-env
source tda-env/bin/activate
pip install gudhi ripser giotto-tda scikit-learn torch matplotlib
\end{verbatim}
\end{warning}

%% ----------------------------------------------------------------
\section*{Historical Milestones}

Topological Data Analysis draws from several foundational contributions:

\begin{itemize}
  \item \textbf{1990s}: Edelsbrunner, Letscher, and Zomorodian introduce
        the first notions of persistent homology.
  \item \textbf{2002}: Zomorodian and Carlsson formalize persistent
        homology and its computation via graded module decomposition.
  \item \textbf{2005}: Cohen-Steiner, Edelsbrunner, and Harer prove
        the fundamental stability theorem for persistence diagrams.
  \item \textbf{2007}: Singh, M\'emoli, and Carlsson introduce the
        Mapper algorithm.
  \item \textbf{2009}: Chazal, Cohen-Steiner, Glisse, Guibas, and Oudot
        extend stability results to Wasserstein distances.
  \item \textbf{2015--present}: Growing integration of TDA with machine
        learning and deep learning.
\end{itemize}

%% ----------------------------------------------------------------
\section*{Acknowledgments}

This course owes much to the pioneering work of Gunnar Carlsson, Herbert
Edelsbrunner, Fr\'ed\'eric Chazal, Steve Oudot, and the entire TDA community.
We also thank the developers of GUDHI, Ripser, and giotto-tda for the
quality of their tools.

\begin{intuition}
TDA rests on a simple but profound idea: \emph{the shape of data carries
information}. The tools of algebraic topology allow us to quantify this
shape in a rigorous, stable, and computable manner. This course will give
you the keys to exploit this idea in your own research.
\end{intuition}

\bigskip
\hfill\textit{The authors, March 2026}
